% Generated draft: requires mathematical and source review.
\documentclass[11pt]{article}
\usepackage[margin=25mm]{geometry}
\usepackage{fontspec}
\setmainfont{Latin Modern Roman}
\usepackage{amsmath,amssymb,mathtools,unicode-math}
\setmathfont{Latin Modern Math}
\usepackage{enumitem,needspace,xurl,hyperref}
\setlength{\parindent}{0pt}
\setlength{\parskip}{5pt}
\setlength{\emergencystretch}{3em}
\newcommand{\N}{\mathbb{N}}
\newcommand{\Z}{\mathbb{Z}}
\newcommand{\Q}{\mathbb{Q}}
\newcommand{\R}{\mathbb{R}}
\newcommand{\C}{\mathbb{C}}
\newcommand{\F}{\mathbb{F}}
\newcommand{\A}{\mathbb{A}}
\newcommand{\PP}{\mathbb{P}}
\newcommand{\E}{\mathbb{E}}
\newcommand{\eps}{\varepsilon}
\newcommand{\dd}{\,\mathrm{d}}
\begin{document}
\section*{UnsolvedMath ID: 30004044}
\textbf{Source problem code:} catalogue entry
\subsection{Definitions and mathematical statement}
Use the usual convention that graphs are finite, simple, and undirected; the supplied record does not explicitly state this convention. For a positive integer \(n\), let \(K_n\) denote the complete graph on \(n\) vertices. A graph is \(r\)-regular if every vertex has degree \(r\). A decomposition of \(K_n\) into copies of \(F_1,\ldots,F_t\) means that there are spanning subgraphs \(H_1,\ldots,H_t\) of \(K_n\) such that
\[
H_i\cong F_i\quad(1\leq i\leq t),\qquad
E(K_n)=\bigsqcup_{i=1}^{t}E(H_i).
\]
The copies may use different identifications of their vertices with \(V(K_n)\), and the graphs \(F_i\) need not be distinct or connected.

\textbf{Question.} Is it true that for every nonnegative integer \(\Delta\) there exists a positive integer \(n_0=n_0(\Delta)\) such that, for every integer \(n\geq n_0\), every positive integer \(t\), and every family of \(n\)-vertex graphs \(F_1,\ldots,F_t\), if each \(F_i\) is \(r_i\)-regular with
\[
r_i\in\Z_{\geq 0},\qquad r_i\leq\Delta\quad(1\leq i\leq t),\qquad
\sum_{i=1}^{t}r_i=n-1,
\]
then \(K_n\) has a decomposition into copies of \(F_1,\ldots,F_t\)?

The integer domains above make explicit the usual interpretation of a degree bound. Zero-degree factors are allowed because the supplied statement does not exclude them. No additional parity restriction on \(n\) is imposed: the hypotheses already require the existence of the specified regular graphs, and therefore require \(nr_i\) to be even for every \(i\). The supplied record labels this question open; this draft does not independently verify that status.
\subsection{Short English statement}
For each fixed degree bound, must every sufficiently large complete graph admit an edge partition into any prescribed family of regular spanning graphs whose degrees obey that bound and sum to \(n-1\)? The intended graph convention is taken to be finite, simple, and undirected, subject to source confirmation.
\subsection{Sources}
\begin{itemize}
\item Local UnsolvedMath registry, permanent ID \(30004044\), source problem code \texttt{catalogue entry}, title: \emph{Universal Decomposition of Complete Graphs into Bounded-Degree Factors}. The record belongs to the Oberwolfach Reports Open Problems collection and gives proposed year 2019; no individual proposer is supplied.
\item Supplied report citation: \emph{Graph Theory (2020)}, Oberwolfach Reports, DOI 10.4171/owr/2019/1, pp.~5--63. \url{https://doi.org/10.4171/owr/2019/1}. The literature list labels this same reference \emph{Oberwolfach Report 1/2019}; the supplied year descriptions are retained without reconciliation.
\item Supplied literature reference: \emph{Resolution of the Oberwolfach problem}. \url{https://web.mat.bham.ac.uk/D.Kuehn/OW.pdf}.
\item Supplied literature reference: \emph{Recent survey on decomposition problems}. \url{https://eprints.lse.ac.uk/126935/1/978-3-98547-564-3-volume-6-chapter-5567.pdf}.
\item Literature assessment dated exactly \texttt{2026-08-22}: the supplied curation labels the problem open and states that known theorems cover important special regimes but not arbitrary bounded-degree factor families. This is supplied dated curation derived from the OWR collection, not an independently verified literature assessment; time-sensitive status and cited claims require independent checking.
\item The registry marks the statement \texttt{corrected\_verified} and says that its mathematical content was preserved while notation, typography, duplicate clauses, or omitted definitions were repaired. These are supplied editorial claims, not mathematical verification performed for this draft.
\end{itemize}
\end{document}
